Zijn jachtgeweer was lang.
Emmy was daardoor bang

dat de spuit bij een vlucht
de afstand overbrugt

en haar zal gaan raken.
Ze moest haar plan staken.

Gesterkt door zijn verhaal
dat hij fenomenaal

scherp ermee kon schieten.
Dat hij moest loodgieten

ook in een medemens
was de dringende wens

van de bisschop destijds.
Dat was niet wederzijds.

Want voor een nederlaag
wilde hij zich niet graag

laten afpeigeren.
Hij kon niet weigeren,

dus liet hij zich testen
om het te verpesten.

Men plaatste een schietschijf
op een meter of vijf.

Hij schoot toch drie keer mis
tot grote ergernis

van de testofficier.
Hij was in het vizier

omdat iedereen zei:
“Hem wil je er graag bij,

want hij schiet alles raak.”
Dat hoorde men zo vaak

dat ze de jager dus
vroegen voor deze klus.

Maar de keurder gaf op.
De keuring bleek een flop.

Nadat het schietterrein
was ontdaan van chagrijn

ging de jager terug.
Hij legde op zijn rug

een konijn, een fazant,
en een vos in een mand,

die zich als betweters
op honderden meters

buiten bereik achtten.
Dat viel te betrachten.

“Ik wil dat hij aandrinkt,
maar niet verder aandringt”,

dacht Emmy na een zucht
met de jagers dranklucht.

“Laat hem zich verdoven.”
Ze wilde geloven

dat hij haar niet versperd
als hij stom dronken werd.

Flink bedwelmd door zijn stank,
wist Emmy niet dat drank

zelfmedicatie was.
Zijn geestespijn genas.

Op de korte termijn
bracht het zeker welzijn.

Het werkte kalmerend,
maar ook stimulerend.

Ooit uit sociale angst
deed alcohol zijn vangst.

In gebruik doorgedraafd
maakte het hem verslaafd.

De drankgeur gaf Emmy
een breinirritatie.

Pijn als een bevalling
die door haar hersens ging,

opende de sluizen
en liet haar verhuizen

naar een ver verleden
in zevenmijls treden.

Felle bliksemschichten
deden haar geest zwichten.

Een wervelwind van licht
bracht het gekwelde wicht

met enorme snelheid
naar een vergane tijd.

Waarbij de waarneemster
niet alleen deelneemster

was maar de reis inhield.
Zich tot de trip verhield

“als zijn tot het wezen”.
Zoals men kan lezen

in het werk van Eckhart.
Het vormt daarnaast het hart

van quantumfysica,
waar de mechanica

van elk quantumdeeltje
slechts een tafereeltje

is uit een verdeling
bepaald na een meting.

Migraines scherpe zwaard
bleek de hemelreis waard.

De les was gekozen
en wel uitgeplozen

voor het fysieke rijk.
Daar was het belangrijk.

Wantrouwen als thema.
De mens als het schema.

De eerste pagina
van iemands samsara,

met meer tussenstops.
Waarin Thomas Hobbes’

wantrouwen als aanleg
de te verkennen weg

was voor Maria’s ziel.
Niet dat het haar beviel.

Zaad dat zijn weg baande
kwam na negen maanden,

nadat het zich loosde,
als eerste geboorte

in de buitenwereld.
De ziel werd gezeteld

in een passend wezen.
God was weer herrezen

vanuit diepe vrede
en liefde alsmede

een niet-fysieke staat
zonder elk goed of kwaad

of licht of duisternis
in de gevangenis

van een ontluikend brein.
Om zijn oneindig zijn

vanuit de beperking
van de veroordeling

middels een aanklager
als de fakkeldrager

verder te verkennen
en te leren kennen.

Al het goddelijke
en het wereldlijke

werden weer verenigd.
Hun brug meer verstevigd.

Het verbond één op één
alle levens aaneen.

Zo op 11:11
ontwaakte er een zelf,

nog niet eerder bevrijd
uit de lege droomtijd,

in een vurige flits.
Begeleid door een gids.

Een magisch wijze elf,
gekoppeld aan dit zelf

tijdens haar levensduur,
was één met de natuur

en haar cyclische aard,
die hij helpt en bewaard.

Deze ervaren ziel
met zijn begaan profiel

steunde met suggestie
deze incarnatie.

Het kind kwam in het licht
als een klein lichtgewicht

met de helm geboren.
Ze was uitverkoren

voor voorspoed en geluk.
Men trok het vlies vlug stuk.

Een heilig feest seinde
“de oogst is ten einde”

op de kille herfstdag.
Dat gold voor de veldslag

bij Viadangos ook.
Men zag nog wel de rook,

maar men lag toch op koers
om Martinus van Tours

waardig te herdenken,
door aandacht te schenken

aan zijn vrome daden.
De mensen aanbaden

hem gedurende de mis.
Bij die gebeurtenis,

als kiem van licht en hoop
met een lange aanloop,

werd Maria verwekt.
En na een lang traject

eindelijk geboren.
Zo kon Emmy horen.

Ze wist dat zij daar lag.
Toen ze de baby zag

zwevend van bovenaf.
Haar ziel die zich vrijgaf.

Angsten voor het leven
werden opgeheven

en maakten plaats voor moed.
Er kwam hoop op voorspoed.

De landstreek beheerste
Alfonso, de wreedste

heerser van Aragón.
Urraca van León,

en ook van Castilla,
was destijds zijn ega

om de beide rijken
tactisch glad te strijken

tegen de gehate
druk van moslimstaten

en interne krachten,
die ze storend achtten.

Hun gefikst huwelijk
bleek al snel gruwelijk.

Hun emoties kotsten.
Hun ambities botsten

als strijders in harnas.
Met “De Fervencias

de Ávila” begon
de heer van Aragón

door 60 gijzelaars,
als vredematelaars

die elk verzet smoorden,
uit wraak te vermoorden,

vanwege hun ontrouw,
de oorlog met zijn vrouw.

De strijd begon pas echt
door een lokaal gevecht

op louter een steenworp
van het Iberisch dorp.

Waar Maria ontstond
op die heilige grond

en zich zou opwerken.
Zou men haar opmerken?
